\begin{table}[htbp]
\centering
\setlength{\tabcolsep}{3.5pt}
\footnotesize
\caption{The predictions of the experimental design and the outcome of each: the four statements of the measurements on the default definition, informative withdrawals and thin grades, the closed forms of the Theory section, and the conjecture of that section, with the quantity that decides each.}
\label{tab:outcomes}
\begin{tabular}{lp{5.2cm}p{4.2cm}p{2.6cm}}
\toprule
Id & Prediction & Result & Outcome \\
\midrule
\relax H1 & \relax On S\&P's CEREP cohorts the transition page's long-run PD of CCC to C is below one half of the default page's (95\% interval of the ratio below 0.5), and the risk weight of the cohort mix falls by more than 5 points & $0.323$ [0.259, 0.391]; $10.5$ points [8.3, 13.1] & \relax held (the point values reproduce the committed ones; the bootstrap interval is new) \\
\relax H2 & \relax With $k=9$ the pooled B PD with withdrawals removed and with withdrawals followed are each at least 10\% below the truth (95\% interval of the relative bias below $-10$\%) & \relax removed $-35.1\pm0.3$, followed $-15.6\pm0.3$ & \relax held \\
\relax H3 & \relax At $k=9$ neither the test of time homogeneity nor the test of the last matrix rejects more than 10\% of the repetitions at the 5\% level (upper Wilson bound below 10\%). The statement behind it, that no transition test of the engine sees the withdrawals, is refuted by a rate whose lower bound is above 10\% at any $k$ or any test & $k=9$: $5.8$ [4.5, 7.4] and $1.4$ [0.8, 2.3]; last matrix at $k=19$: $16.3$ [14.1, 18.7]; Markov order at $k=9$: $8.3$ [6.7, 10.2] & \relax held for the two tests at $k=9$; the statement is not held, the last-matrix test sees the withdrawals at $k=19$ \\
\relax H4 & \relax With 50 obligors the Wald interval covers the PD of AAA in $0.33\%$ of the samples (to two digits) and the Jeffreys and Agresti-Coull intervals cover at least 95\% in every grade & \relax Wald $0.33$\%; minimum $96.0$\% (Jeffreys), $95.5$\% (A-C) & \relax reproduced (an exact enumeration of values already committed; it could not fail) \\
\relax H5a & \relax The closed forms of Proposition~\ref{prop:first-year} agree with the simulated means of the first cohort within three Monte Carlo standard errors & \relax largest $|z|=1.69$ over 30 distinct comparisons & \relax held \\
\relax H5b & \relax The numerical evaluation of the model's masses over five cohorts, with the horizon truncation, agrees with the pooled simulated means within three Monte Carlo standard errors & \relax largest $|z|=1.59$ over 40 distinct comparisons & \relax held \\
\relax Conj. & \relax The first-order formula of Equation~\eqref{eq:first-order} describes the first-cohort bias of a grade with a small PD to within three standard errors & \relax CCC to C, removal: $z=3.02$, $3.31$, $6.32$ at $k=2,5,19$; B, keeping, $k=19$: $z=-3.07$ & \relax refuted for CCC to C and, at $k=19$, for keeping in B \\
\bottomrule
\end{tabular}
\end{table}
